
\documentclass{article}
\usepackage{colortbl}
\usepackage{makecell}
\usepackage{multirow}
\usepackage{supertabular}

\begin{document}

\newcounter{utterance}

\centering \large Interaction Transcript for game `guesswhat', experiment `Level\_3', episode 2 with Qwen3{-}30B{-}A3B{-}t0.0.
\vspace{24pt}

{ \footnotesize  \setcounter{utterance}{1}
\setlength{\tabcolsep}{0pt}
\begin{supertabular}{c@{$\;$}|p{.15\linewidth}@{}p{.15\linewidth}p{.15\linewidth}p{.15\linewidth}p{.15\linewidth}p{.15\linewidth}}
    \# & $\;$A & \multicolumn{4}{c}{Game Master} & $\;\:$B\\
    \hline

    \theutterance \stepcounter{utterance}  
    & & \multicolumn{4}{p{0.6\linewidth}}{
        \cellcolor[rgb]{0.9,0.9,0.9}{
            \makecell[{{p{\linewidth}}}]{
                \texttt{\tiny{[P1$\langle$GM]}}
                \texttt{Let's play a guessing game. Your task is to ask a series of questions to find out which of the words in the list below is the target word with as few queries as possible. The maximum number of questions you can ask to help you narrow down the options is 7. Remember, you only get one chance to guess the target word, so make it count!} \\
\\ 
\texttt{The rules of the game are: } \\
\\ 
\texttt{1. Only ask one question per turn. DO NOT ASK MORE THAN ONE QUESTION.} \\
\\ 
\texttt{2. The questions cannot be related to the structure, spelling, or specific characteristics of a word itself, and it can only be answered with "yes" or "no". } \\
\\ 
\texttt{3. Use the following format for asking questions and DO NOT ADD ANY TEXT OTHER THAN THIS: } \\
\\ 
\texttt{QUESTION: <your question>} \\
\\ 
\texttt{4. You have only one opportunity to guess the target word. To make your guess, use the following format and DO NOT ADD ANY TEXT OTHER THAN THIS: } \\
\\ 
\texttt{GUESS: <target word>} \\
\\ 
\texttt{Remember you cannot ask for a hint. Make your question brief. Let’s start!} \\
\\ 
\texttt{List of candidate words: {[}'Nephew', 'Uncle', 'Grandfather', 'Grandmother', 'Father', 'Mother', 'Wife', 'Husband'{]}} \\
            }
        }
    }
    & \\ \\

    \theutterance \stepcounter{utterance}  
    & \multicolumn{4}{p{0.6\linewidth}}{
        \cellcolor[rgb]{0.9,0.9,0.9}{
            \makecell[{{p{\linewidth}}}]{
                \texttt{\tiny{[P1$\rangle$GM]}}
                \texttt{<think>} \\
\texttt{Okay, let's see. I need to figure out the target word from the list using as few yes/no questions as possible, max 7. The words are all family relations: Nephew, Uncle, Grandfather, Grandmother, Father, Mother, Wife, Husband.} \\
\\ 
\texttt{First, I should think about the possible categories or attributes that can differentiate them. Since I can't ask about spelling or structure, I need to think about their meanings. Maybe gender? But some words are gender{-}neutral, like 'Parent' (but here it's Father/Mother). Wait, the list has Father, Mother, which are gendered. But others like Uncle, Aunt (but Aunt isn't here), Nephew, etc. Wait, the list includes 'Wife' and 'Husband', which are gendered. So maybe a question about gender? But how to phrase it without mentioning the word itself. Like, "Is the target word a male relative?" But then, some words are male (Uncle, Nephew, Father, Husband) and others female (Grandmother, Mother, Wife). But Grandfather is male, Grandmother female. So if I ask "Is the target word a male?" then the answer would be yes for Uncle, Nephew, Father, Husband, Grandfather. No for the others. But then I can split the list into two groups. But maybe that's a good first question.} \\
\\ 
\texttt{Alternatively, maybe ask about the} \\
            }
        }
    }
    & & \\ \\

    \theutterance \stepcounter{utterance}  
    & & & \multicolumn{2}{p{0.3\linewidth}}{
        \cellcolor[rgb]{0.9,0.9,0.9}{
            \makecell[{{p{\linewidth}}}]{
                \texttt{\tiny{[GM$|$GM]}}
                \texttt{Invalid format. Guesser must use the form 'QUESTION: ' or 'GUESS: '.} \\
            }
        }
    }
    & & \\ \\

    \theutterance \stepcounter{utterance}  
    & & & \multicolumn{2}{p{0.3\linewidth}}{
        \cellcolor[rgb]{0.9,0.9,0.9}{
            \makecell[{{p{\linewidth}}}]{
                \texttt{\tiny{[GM$|$GM]}}
                \texttt{abort game} \\
            }
        }
    }
    & & \\ \\

\end{supertabular}
}

\end{document}
